% Shared preamble for the Radioactive Signal Explorer explainers (docs/signals/*.tex).
\documentclass[11pt]{article}
\usepackage[letterpaper,margin=1in]{geometry}
\usepackage[T1]{fontenc}
\usepackage{lmodern}
\usepackage{microtype}
\usepackage{amsmath,amssymb,mathtools,bm}
\usepackage{siunitx}
\usepackage{graphicx}
\usepackage{booktabs,array,tabularx}
\usepackage{xcolor}
\usepackage{enumitem}
\usepackage{caption}
\usepackage[section]{placeins}
\usepackage[most]{tcolorbox}
\usepackage[colorlinks=true,linkcolor=blue!50!black,citecolor=blue!50!black,urlcolor=blue!50!black]{hyperref}
\usepackage[capitalise,noabbrev]{cleveref}
\sisetup{group-separator={,},group-minimum-digits=4,separate-uncertainty}
\definecolor{labblue}{HTML}{0088EE}
\definecolor{exbg}{HTML}{FFF7EC}
\definecolor{boxbg}{HTML}{EEF4FB}
\newtcolorbox{keyidea}[1][]{colback=boxbg,colframe=labblue,boxrule=0.6pt,arc=2pt,left=6pt,right=6pt,top=4pt,bottom=4pt,fonttitle=\bfseries,title={#1}}
\newtcolorbox{pitfall}[1][]{colback=red!3,colframe=red!55!black,boxrule=0.6pt,arc=2pt,left=6pt,right=6pt,top=4pt,bottom=4pt,fonttitle=\bfseries,title={#1}}
\newcounter{exercise}
\newtcolorbox{exercise}[1][]{colback=exbg,colframe=orange!70!black,boxrule=0.6pt,arc=2pt,left=6pt,right=6pt,top=4pt,bottom=4pt,fonttitle=\bfseries,
  title={Exercise~\refstepcounter{exercise}\theexercise\ifx\relax#1\relax\else\ (#1)\fi}}
\newcommand{\dd}{\mathrm{d}}
\newcommand{\keV}{\,\mathrm{keV}}
\newcommand{\MeV}{\,\mathrm{MeV}}
\newcommand{\me}{m_e c^2}
\newcommand{\nuc}[2]{${}^{#1}$#2}
\newcommand{\code}[1]{\texttt{#1}}
\newcommand{\explorer}{\href{https://dark-matter.ucr.edu/Signals.html}{Radioactive Signal Explorer}}
\newcommand{\seriesnote}{\noindent\small This note is one of five explaining the physics behind the \explorer: gamma rays, alpha decay, $(\alpha,n)$ neutrons, beta decay and Cherenkov light. Every figure was computed with the page's own code (\code{assets/signals-physics.js}); the \LaTeX{} sources and figure scripts are in \code{docs/signals/} of \url{https://github.com/UCRDarkMatter/dark-matter-lab-site}. Exercises are marked (A) analytic, (P) programming, (R) open-ended.\normalsize}
